\section{Experimental Results} \label{sec:experimental_results}

This section presents the results of the experiments conducted with the implemented autonomous inspection process for a rotating wind turbine. Firstly, the details on the experiment setting will be introduced, followed by the results obtained through the experiments. Lastly, the section will end with a small discussion on the matter.

% data on the software base
The work done in this study relies heavily on the code bases for autonomous UAV operation and point cloud registration and inspection developed by the LARICS \footnote{Laboratory for Robotics and Intelligent Control Systems - \hyperlink{https://larics.fer.hr/}{larics.fer.hr}} laboratory from the University of Zagreb, similar to the work done in \cite{ref:Car2020}. The autonomous UAV operation code base takes care of the low level control providing an interface for higher level navigation and inspection scheduling. On the other hand, the point cloud registration and inspection code base provides high level handling of point cloud data coming from the LiDAR sensors, including operations such as registration, ground plane filtering, clustering, and matching.

% data on the uav with which the bags were recorded
The data used to perform experiments was collected by a custom UAV equipped with an array of sensors. The system runs on an Intel NUC with ROS as the communication base. Moreover, the flight control is done through the Pixhawk and Ardupilot firmware. The positioning is done through a fusion of GPS and IMU signals, and based on the run, the UAV is fitted with either a Velodyne or an Ouster LiDAR sensor.

\begin{table}[H]
    \centering
    \begin{tabular}{||c|c||}
    \hline
     \multicolumn{2}{||c||}{Pometeno Brdo wind farm} \\
     \hline
    \hline
        Location & near Split, Croatia \\
    \hline
        Base height & 45 m \\
    \hline
        Rotor radius & 30 m \\
    \hline
        Turbine number & 17 \\
    \hline
        Electricity generation & 84,7 MW \\
    \hline
    \end{tabular}
    \caption{Inspected wind turbine information}
    \label{tab:turbine_info}
\end{table}

% data on the setting in which the bags were recorded
The data collection process was done on a wind turbine farm near Split, Croatia with details given in~\ref{tab:turbine_info}. It involved bringing the UAV up in front of a wind turbine and hovering in front of it for an amount of time, all the while collecting information about the environment with the LiDAR. The data is stored ROS .bag files for playback and further experiments in simulation. This data was used to validate the process of localization of the rotating turbine with the algorithm implemented in this study.



\begin{table}[H]
    \centering
    \begin{tabularx}{0.48\textwidth}{||X|X|X|X||}
    \hline
     Run no. & LiDAR & Match results & Yaw (rad) \\
     \hline
    \hline
        1 & Ouster & 0.262016 & 1.73169\\
    \hline
        2 & Ouster & 0.251234 & 1.98284\\
    \hline
        3 & Ouster & 0.277545 & 2.187747\\
    \hline
    \end{tabularx}
    \caption{Experimental results}
    \label{tab:results}
\end{table}


The results of the experiments are given in table~\ref{tab:results}. The experiments are performed with the LiDAR sensor Ouster. The principal dimensionless metric for the matching process is the match score between one of the clusters in the registered point cloud and the 3d model presenting the wind turbine. The value of the match score decreases with higher match quality. Furthermore, the yaw value representing the yaw angle of the matched wind turbine in the scene. Generally, for clusters not resembling the wind turbine the match score is in the 100-200 range, for clusters similar to a part of the turbine 1-100, and for a good match a score <1 is expected. \\
Let us discuss the matching results for the Ouster LiDAR (the first three rows in the table). These represent experiments done on January 25th 2022 and as can be seen from the match results, all of them successfully matched the 3d model onto the registered point cloud. The score of all three experiments is in the 0.25-0.28 range, resulting in an almost perfect fit. This can be further confirmed visually with matches similar to~\ref{fig:matching_4}. Moreover, from yaw values being different for every run, it is clear that the system was initialized at a different points with regards to the wind turbine.\\ 

%October 28th Ouster